\documentclass[
		a4paper,
		12pt,
		twoside,
		openany,        % Chapter beginnen auf rechter Seite
		parskip,
    % tocbasic,
		final,          % DEBUG modus
		11pt,
    % draft, # DEBUG MOUS
		appendixprefix,
    % headsepline,                       % Linie nach Kopfzeile.
    % footsepline,                       % Linie vor Fusszeile.
    % bibtotocnumbered,                  % Literaturverzeichnis im Inhaltsverzeichnis nummeriert einfügen.
		chapterprefix,  % Kapitel anschreiben als Kapitel
]{scrreprt}
\usepackage[twoside,width=14cm, height=24cm,right=3.25cm, left=2.5cm]{geometry}
\setlength{\marginparwidth}{2.25cm}
\clubpenalty = 10000
\widowpenalty = 10000 \displaywidowpenalty = 10000

\usepackage[ngerman]{babel}

\usepackage[utf8]{inputenc}		%Damit kann man auch Umlaute direkt eintippen  \o/
\usepackage{textcomp}
\usepackage{amsmath}
\usepackage{amssymb}
\usepackage{amsbsy}
\usepackage{amsfonts}
\usepackage{amsthm}
\usepackage{graphicx}	%including images
\usepackage{fancyhdr}	%for editiable headers and footers
\usepackage{setspace} 	%for 1,5times row space
\usepackage{color}
\definecolor{darkblue}{rgb}{0,0,.5}
\usepackage[colorlinks,
  pdfpagelabels,pdfstartview=FitH,bookmarksopen=true,bookmarksnumbered=true,
  linkcolor= black,
  plainpages=false,
  hypertexnames=false,
  citecolor=black,linkcolor=black,menucolor=black,urlcolor=darkblue]
  {hyperref}
			%clickable TOC
\usepackage{multirow}	%linebreak in Table
\usepackage{enumerate}	% \begin{enumerate}[a.]     Custom Nummerierung
\usepackage{float}	% Images mit Float
\usepackage{colortbl}	%colorTables
\usepackage{framed}   		%begin{framed}
\usepackage{longtable} 		%begin{longtable}
\usepackage{makeidx}
\usepackage[square,sort&compress,comma,numbers]{natbib}  % natbib, aber Zitiert mit [..] anstatt (..)
\usepackage{epigraph}		% Epigraph = stylische Zitatte
\usepackage{listings}	% Quellcode Listings und List of Listings
\usepackage{caption}	% Modifizierung und setzen von custom captions
\usepackage[dvipsnames,table]{xcolor}
\usepackage[nottoc]{tocbibind}	% TOC Plugin -> ??
\usepackage{relsize}
\usepackage[toc]{glossaries}
%\usepackage{marginnote}
%\usepackage{minted}  % TODO Pygments wrapper... ohje :)



%%%%%%%%%%%%%%%%%%%%%%%%%%%%%%%5
%% Metainfos
%%%%%%%%%%%%%%%%%%%%%%%%%%%%%%%

\newcommand{\ownTitle}{Diplomarbeit}
\newcommand{\ownTitleZ}{Testgetriebene Entwicklung von Web-Serveranwendungen auf der Basis von Ruby on Rails}
\newcommand{\ownAutor}{Stefan Wienert, HTW Dresden}

\title{\ownTitle}
\author{\ownAutor}
\pdfinfo{
 /Title  (\ownTitle)
 /Author (\ownAutor)
 /Subject (\ownTitleZ)
 /Keywords()
}
\date{\today}
\setlength{\epigraphwidth}{.8\textwidth}


%%%%%%%%%%%%%%%%%%%%%%%%%%%%%%%%%%
%% HEAD AND FOOTLINES
%%%%%%%%%%%%%%%%%%%%%%%%%%%%%%%%%%
\pagestyle{headings}
\definecolor{Gray}{rgb}{0.85,0.85,0.90}

\makeindex 	%start Aufzeichnung eines Stichwortverzeichnisses
\makeglossary  	%start Aufzeichnung eines Glossars

%%%%%%%%%%%%%%%%%%%%%% STYLINGS & FONTS %%%%%%%%%%%%%%%%%%%%%%%%%%%%
\renewcommand{\footrulewidth}{0.5pt}
%\let\cleardoublepage\clearpage
\parskip 5pt       % sets spacing between paragraphs
\parindent 0pt


\usepackage{bookman}  % 118 pages, sehr sauber
% \usepackage{newcent}  % 116 pages, nice aber etwas zu kraeftig
% \usepackage{lmodern}  % 109 pages, standardmaessig
% \usepackage{boisik}    % 113 pages, gewagt aber sehr detailiert, etwas ungerade
% \usepackage{palatino}    % 111 pages, Arial'y sehr langweilig in den Ueberschriften
% \usepackage{cmbright}    % 111 pages, Wie ein modernes Buch. Gut fuer ne Doku oder so
 \usepackage{charter}    % 109 pages, Wie ein modernes Buch. Gut fuer ne Doku oder so
%\fontsize{12}{15}

\usepackage{hyphenat}
\usepackage[T1]{fontenc}


\hyphenation{
Com-mu-ni-ty
Soft-ware-Ent-wick-lung
Soft-ware
Aus-für-lich
Um-for-mu-lie-rung
Test-fall-ab-deck-ung
Test-Unit
Test/Unit
Da-ten-bank-an-bin-dung
Tes-ten
Grob-ent-wurf
Pro-gram-ming
Active-Record
Ini-ti-ali-sie-rung
Sys-tem
be-schreibt
De-ckungs-gleich-heit
Ba-sis-at-tri-bu-te
Ar-gu-ment-Er-ror
}

\usepackage{setspace}
\setstretch{1,15} % Mein TOC ist zu lang!!
\newcommand{\fullref}[1]{\ref{#1} -- \glqq\,\textit{\nameref{#1}}\grqq}
% \fullref{sec:bla}  -> "8.2 - 'Name der Subsection"
% zum Verweis mit Wiederholung des Verweistitels
\newcommand{\fullref}[1]{\ref{#1} -- \glqq\,\textit{\nameref{#1}}\grqq}


% graphik on the right or left border (dep. on twosided)
\newcommand{\bordergraphic}[1]{\marginline{\includegraphics[width=0.8\marginparwidth]{#1}}}

% text on the border
\newcommand{\randbem}[1]{\marginpar[\fontsize{9}{10}\selectfont \raggedleft{#1}]
  {\fontsize{9}{10}\selectfont \raggedright{#1}}
}
% quote on the border
\newcommand{\borderquote}[2]{\setlength{\epigraphwidth}{\marginparwidth}
\marginpar{\fontsize{9}{10} \epigraph{#1}{#2}}
\setlength{\epigraphwidth}{0.8\textwidth}
}


% Einheitliche Bildquelle als zusaetzliches Label
\newcommand{\imgsource}[1]{\captionsetup{font={footnotesize,bf,it}} \caption*{#1}}

% convenient long arrow
\newcommand{\arr}{$ \Longrightarrow$}


% load ruby as a standard listing
\lstloadlanguages{Ruby}
\lstset{%
  basicstyle=\ttfamily\footnotesize\color{black},
  commentstyle = \ttfamily\color{gray},
  keywordstyle=\ttfamily\color{blue},
  stringstyle=\color{red},
  showspaces=false,               % show spaces adding particular underscores
  showstringspaces=false,
  frame=single,
  breaklines=true
}
% Allow umlauts in lstlistings
\lstset{literate=%
{Ö}{{\"O}}1
{Ä}{{\"A}}1
{Ü}{{\"U}}1
{ß}{{\ss}}2
{ü}{{\"u}}1
{ä}{{\"a}}1
{ö}{{\"o}}1
}

% fancy glossar reference with arrow in front of
\newcommand{\glossar}[1]{$^\uparrow$\gls{#1}}
\newcommand{\glossarpl}[1]{$^\uparrow$\glspl{#1}}

  % make table beautiful:
  %   zebra colors, more difference, header bg color
  % cannot be earlier, because frontpage uses tables for positioning
\newcommand{\fancytables}{\rowcolors{1}{lightgray!50!white}{white}
\colorlet{tableheadcolor}{gray!40}
\renewcommand\arraystretch{1.25}% (MyValue=1.0 is for standard spacing)
}

%\newglossaryentry{test} {
  name={Testfall},
  description= { oder \index{Test} ist eine, meist automatisierte, Prüfung des Programmverhaltens bei definierten Eingabeparametern},
  plural={Tests}
}

\newglossaryentry{testrunner} {
  name={Test-Runner},
  description= {oder Testtreiber) ist eine Software, um \glossarpl{test} aufzurufen und deren Ausführung zu überwachen und zu steuern}
}
\newglossaryentry{rake} {
  name={Rake},
  description= {ist ein Build-Programm, das im Gegensatz zu \texttt{make} mit Ruby programmiert wird. Es ist modular mit eigenen Tasks erweiterbar. Innerhalb von Rails dient es auch als weitere Schnittstelle zur Anwendung, um Wartungsaufgaben auszuführen (z.B. Datenbankbackup oder Cronjobs)}
}
\newglossaryentry{racktest} {
  name={RackTest},
  description= { ist eine Test-API für Rack-Anwendungen. Rack ist ein minimales Interface zur Kommunikation mit Webservern, Middlewares und Webframeworks, wie z.B. Rails. RackTest kommuniziert direkt mit dem Rack-Interface, verzichtet auf HTTP Traffic und ist damit sehr schnell}
}

\newglossaryentry{testsuite} {
  name={Test-Suite},
  description= {ist eine Gruppe von mehreren Testfällen für eine Komponente oder für das gesamte System}
}

\newglossaryentry{fehler} {
  name={Software-Fehler},
  description= {ist jede Abweichung der Implementierung von der Spezifikation \citep[IEEE/ANSI]{fritzsche_sqm_2005}. Diese Abweichungen können zu fehlerhaftem Verhalten der Anwendung führen},
  text=Fehler,
  first={Software-Fehler}
}
\newglossaryentry{quality} {
  name={Code-Qualität},
  description= {beinhaltet die Software-Qualitäten Lesbarkeit, Testbarkeit, Wartbarkeit, Erweiterbarkeit, geringe Komplexität},
}

\newglossaryentry{metrik} {
  name={Code-Metrik},
  plural={Code-Metriken},
  description={Eine Softwaremetrik ist das Ergebnis einer statischen oder dynamischen Codeanalyse zur Generierung von Informationen über den Quelltext. Beispiele: Testabdeckung, Anzahl Codezeilen, Anzahl \glossarpl{smell} pro Codezeile}
}
\newglossaryentry{refaktorisieren} {
  name = {Refaktorisieren},
  description = {Ist eine Modifikation des Programmcodes ohne Modifikation des externen Verhaltens um nicht-funktionale Eigenschaften des Quellcodes zu verbessern, wie z.B. Lesbarkeit, Wiederverwendbarkeit und Wartbarkeit}
}
\newglossaryentry{Testabdeckung} {
  name=Testabdeckung,
  description={auch: \index{Test!Testabdeckung}Testausführungsabdeckung, Überdeckungsgrad, Testfallabdeckung. Eine dynamische Code-Metrik\index{Code-Metrik}, die angibt, welche Codezeilen durch keinen Test abgedeckt wurden. Es wird unterschieden in die Stufen C0, C1 und C2 mit steigender Komplexität der Messung. Details im Abschnitt \ref{sec:metrics}},
  first = {Testabdeckung}
}
\newglossaryentry{TDD}{
  name={Testgetriebene Entwicklung},
  description={englisch Test Driven Development/Test Driven Design, Ausführlich dargelegt in Abschnitt \ref{sec:tdd}},
  text = TDD,
  first = {Testgetriebene Entwicklung}
}
\newglossaryentry{BDD}{
  name={ Behavior Driven Development},
  description={ (Verhaltensgetriebene Entwicklung). Umformulierung von TDD\index{TDD} zur Ausrichtung auf Businessprozesse. Das Vokabular zielt auf die Spezifikation von Erwartungen im Systemverhalten, anstatt Definition nachträglicher Tests}
}
\newglossaryentry{CRUD} {
  name={CRUD},
  description={Create, Read, Update, Delete sind die vier Basisoperationen, die auf persistente Speicherung ausgeführt werden können},
  first={Create Read Update Delete (CRUD)}
}
\newglossaryentry{smell}{
  name={Code Smell},
  description={oder Bad Smell\index{Code-Smell}. Ist ein Anzeichen für eine suboptimale Quelltextstelle, die auch ein Hinweis auf ein größeres Designproblem sein kann. Oft auch ein Kandidat für eine Refaktorisierung. Informationen zu Smells und deren Refaktorisierung\index{Refaktorisierung} sind im Buch von M.Fowler zu finden \citep{fowler_refactoring_1999}},
  plural={Code Smells}
}
\newglossaryentry{ORM}{
  name={ORM},
  description={Objektrelationales Mapping, ist eine Persistenztechnik, um Objekte transparent in einer Datenbank\index{Datenbank} zu speichern und umgekehrt, Tabellenzeilen in Objekte wieder zu rückübersetzen},
  plural={ORMs},
  first={Objekt-relationales Mapping (ORM)}
}
\newglossaryentry{MVC}{
  name={MVC},
  description={Model-View-Control, ist ein \glossar{patterns}, das insbesondere bei GUI- und Web-Anwendungen beliebt ist. Rails\index{Ruby-on-Rails} basiert auf dem MVC-Muster},
  plural={MVCs},
  first={Model-View-Controller (MVC)}
}

\newglossaryentry{rails}{
  name={Ruby-on-Rails},
  description={ist ein auf der Programmiersprache Ruby basierendes Web-Framework und ist Gegenstand des Kapitels \ref{sec:rails}},
  text = Rails,
  first = Ruby on Rails
}

\newglossaryentry{Test-Double} {
  name={Test-Double},
  plural={Test-Doubles},
  description = {haben die Aufgabe, komplexe Objekte in einem isolierten Test \index{Test}zu simulieren, in dem statt komplexer Berechnungen oder externer Zugriffe konstante Werte geliefert werden. Vertreter dieser Test Doubles sind die Mock\index{Test-Double!Mock}s und Stub\index{Test-Double!Stub}s, siehe Abschnitt
 \ref{sec:mocks}}
}
\newglossaryentry{patterns}{
  name={Entwurfsmuster},
  description={ oder Design Patterns sind bewährte Vorlagen, um häufig wiederkehrende Probleme zu lösen. Weitere Informationen im Buch Design Patterns: Elements of Reusable Object-Oriented Software von Gamma, Helm, Johnson, Vlissides (Gang of Four)}
}
\newglossaryentry{RSS}{
  name={RSS},
  description={  RDF Site Summary, ist ein standardisierter XML-Dialekt zur maschinenlesbaren Verteilung und Veröffentlichung von Inhalten. Es existiert in den Versionen 0.9 bis 2.0.1, die sich nur in Details, wie z.B. Einbindung von Rich-Media (Podcasts, ..) und Namespaces unterscheiden}
}
\newglossaryentry{Test-Umgebung} {
  name={Test-Umgebung},
  description={ist eine spezielle Konfiguration der Hard- und Software, um Tests unter kontrollierten und bekannten Bedingungen auszuführen. Dies beinhaltet neben dem zu testenden Objekten (Object under Test) auch sogenannte \glossarpl{Test-Double}\index{Test-Double} und den \glossar{testrunner}}
}
\newglossaryentry{gem} {
  name={Gem},
  plural={Gems},
  description={als Gem\index{Gem}s (zu Deutsch: Edelsteine) werden in der Ruby-Gemeinschaft Bibliotheken Dritter bezeichnet, die in einem zentralen Repository lagern. Das Bekannteste dieser Repositorien ist rubygems.org. Nahezu alle Ruby-Bibliotheken lassen sich hier finden und innerhalb von Sekunden mittels des Kommandozeilenwerkzeuges gem installieren. Beispielsweise ließe sich \glossar{rails} mittels \texttt{gem install rails} installieren}
}

\newglossaryentry{DSL} {
  name={DSL},
  plural={DSLs},
  description={Eine Domain-spezifische Sprache ist eine auf eine spezielle Problemdomäne ausgerichtete Programmier- oder Spezifikationssprache. Vertreter sind z.B. SQL, die Statistiksprache R oder die Hardwarebeschreibungssprache VHDL. Sie stehen damit den Allzweck-Programmiersprachen gegenüber},
  first={Domain-Specific-Language (DSL)}
}


\newglossaryentry{CMS} {
  name={CMS},
  plural={CMSe},
  description={Content Management System (Inhaltsverwaltungssystem) ist eine, meist Web-basierte Software, die es Nutzern ermöglicht einfach statische Inhalte anzulegen und zu bearbeiten},
first = {Content Management System}
}
% \glsaddall % Enable all glossar entries, even if not used yet

\renewcommand{\glossarypreamble}{Im Folgenden werden einige oft verwendete Begriffe näher erläutert. Innerhalb des Hauptteils dieser Arbeit sind diese Begriffe mit einem $^\uparrow$ gekennzeichnet.}
% \renewcommand{\glossarymark}{Glossar}


  % Glossar definitions

\begin{document}
 \thispagestyle{empty}
\begin{center}
\begin{tabular}{lcr}
 \Large{HTW Dresden} & \verb|       |& \multirow{3}{*}{\includegraphics[height=1.353cm]{material/htwlogo.jpg}} \\
 \Large{Fakultät Informatik/Mathematik} &  & \\
\ownTitle &  & \\
\end{tabular}\end{center}

\begin{verbatim}






\end{verbatim}
\begin{center}
\begin{huge}\textbf{\ownTitle}\end{huge}

\end{center}
\begin{verbatim}



\end{verbatim}
\begin{center}
\textbf{\LARGE{\ownTitleZ}}
\end{center}
\begin{verbatim}




\end{verbatim}
\begin{flushleft}
\begin{tabular}{lll}
\textbf{Autor:} & & Wienert, Stefan\\
& & \\
\textbf{Studiengruppe:} & & 07/041/02\\
& & \\
& & \\
\textbf{Betreuender Professor:} & & Prof. Dr. Fritzsche\\
& & \\
& & \\
\textbf{Datum:} & & \today\\
\end{tabular}

\end{flushleft}
 \cleardoublepage

 \pagenumbering{Roman}
 \tableofcontents		%Inhaltsverzeichnis
 \newpage
 \setstretch{1,25}  % Entspricht \onehalfspacing
 \chapter*{Danksagung}

Besonderen Dank gilt meinem Betreuer, Prof. Fritzsche, der sich regelmäßig und ausführlich mit meiner Arbeit beschäftigte und mir wertvolle Hinweise erteilte.

Mein Dank gilt auch der Fakultät Informatik und Prof. Wiedemann für das zur Verfügung gestellte Büro, der Bibliothek der HTW-Dresden und Prof. Nestler für das Bestellen von tagesaktueller Literatur.

Ich danke auch meiner Frau, die mich während des Schreibens unterstützt hat und bei der Erstellung der Grafiken Tipps gab. Für die orthografische und expressive Optimierung danke ich meinen Freunden und Kollegen Daniel Schäufele, Jörg Sachse, Martin Schneider, Akos Toth und Stefan Koch.

Diese Arbeit entstand in Zusammenarbeit mit der Firma pludoni GmbH und unter Aufsicht des Geschäftsführers Dr. Jörg Klukas. Ich danke ihm für die Möglichkeit, frei zu arbeiten und viele neue Dinge auszuprobieren und für das Vertrauen, das er dabei in mich gesetzt hat.



  % Table -- Zebrastreifen und 1.25 Row Abstand
  \rowcolors{1}{lightgray!50!white}{white}
  \colorlet{tableheadcolor}{gray!40}
  \renewcommand\arraystretch{1.25}% (MyValue=1.0 is for standard spacing)


  \thispagestyle{empty} \mbox{}
  %\newpage
  %\setcounter{page}{-1}
  %\clearpage
  \cleardoublepage
  \pagenumbering{arabic}
 \chapter{Einleitung}
\chapter{Einleitung}

Thema

\section{Motivation}



%\section{Projektbeschreibung und Projektziele}
Für den praktischen Teil dieser Arbeit soll anhand der Entwicklung einer Ruby-on-Rails\hyp{}Anwendung die Testgetriebene\index{TDD} Entwicklung erprobt und angewendet werden.

Als Ergänzung zu den lokalspezifischen Communitys mit strenger Mitgliederauswahl soll nun ein neues, allgemeines IT-Jobportal entwickelt werden. Der voraussichtliche Name wird IT-Jobs\index{IT-Jobs-Projekt}-Und-Stellen.de\footnote{\url{http://www.it-jobs-und-stellen.de/}} sein.

Ziel soll es sein, den regionalen Organisationen auch eine Alternative zu den Community-Mitgliedschaften anzubieten, um kurzfristigen Personalbedarf zu decken, analog zu den bekannten Online-Stellenbörsen stepstone.de, monster.de und jobscout24\footnote{\url{http://www.stepstone.de} - \url{http://www.monster.de} - \url{http://www.jobscout24.de}}. Parallel zu dieser Arbeit wurden die Anforderungen analysiert und die in Abbildung \ref{fig:usecases} dargestellten Anwendungensfälle bestimmt.


\begin{figure}[htbp]
 \centering
 \includegraphics[width=1\textwidth]{./material/usecases.png}
 % usecases.png: 930x899 pixel, 51dpi, 46.50x44.95 cm, bb=
 \caption{Anwendungsfälle}
 \label{fig:usecases}
\end{figure}


\subsection{Anwendungsfälle}
Es gibt zwei verschiedene Nutzertypen. Zum einen der \textit{Bewerber}, der das Ziel hat, einen Job zu suchen, zu finden und sich darauf zu bewerben. Zum anderen der Mitarbeiter eines Kundenunternehmenes (\textit{Arbeitgeber}), der Stellenanzeigen auf der Webseite schalten möchte.
Verbal formuliert sind die Anwendungsfälle die Folgenden:
\begin{itemize}
 \item Ein Bewerber kann die Webseite nach sichtbaren Stellenangeboten durchsuchen.
 \item Ein Bewerber kann die Detailanzeige einer Stelle betrachten.
 \item Ein Bewerber kann sich auf eine Stelle über das System bewerben. Dabei kann er auch eine Verbindung zu seinem Facebook- oder LinkedIn-Account herstellen, falls er das möchte, um so automatisch einen Lebenslauf generieren zu lassen und Stammdaten eintragen zu lassen.
 \item Ein Kunde kann sich am Portal registrieren, dann anmelden und wird so zum \textit{Arbeitgeber}.
 \item Ein Arbeitgeber kann eine Stellenanzeige für ein Vielfaches von 30 Tagen schalten. Damit wird ein Bezahlvorgang über ein Drittanbieterportal (z.B. PayPal) ausgelöst und bei Bestätigung die Stelle sichtbar geschaltet. Er kann dazu zwischen 5 verschiedenen Designs wählen sowie Logo und Bild hinzufügen.
 \item Ein Arbeitgeber kann eine Stelle verlängern oder kopieren, um als sie als Vorlage für ein neues Stellenangebot zu nutzen.
 \item Ein Arbeitgeber kann eine laufende Stelle erneuern (refresh), um so seine Platzierung verbessern. Dies ist möglich, weil die Aktualität einer Stelle Auswirkungen auf die Platzierung in den Suchergebnissen hat. Ein Refresh kostet Geld und löst somit wieder einen Bezahlvorgang aus.
 \item Ein Arbeitgeber kann Statistiken über seine geschalteten Stellenanzeigen einsehen. Dies beinhaltet die Anzahl der Zugriffe pro Tag und Anzahl der Bewerbungen je Stelle.
 \item Ein Arbeitgeber kann seine Stellen mittels eines XML-Feed-Imports automatisiert einlesen. Dazu bezahlt er einen Pauschalbetrag je Kalenderjahr. Das System liest alle 6 Stunden automatisiert diesen Feed ein, der durch HTTP erreichbar sein muss und überträgt gültige Stellenanzeigen in die Datenbank\index{Datenbank}. Nicht mehr in dem Feed enthaltene Stellenanzeigen des Arbeitgebers werden in der Datenbank auf "`unsichtbar"' gesetzt. Das Format ist eine Modifikation von \glossar{RSS} 2.0.
\end{itemize}



\subsection{Nichtfunktionale Anforderungen}
Folgende Rahmenbedingungen und zusätzliche Ziele wurden vereinbart.

\begin{itemize}
 \item Eine hohe C0-\glossar{Testabdeckung} von mindestens 95\% als Grundlage für weitere Arbeiten und für den Prozess der Testgetriebenen Entwicklung\index{TDD}
 \item Eine hohe Erweiterbarkeit, um langfristig auch die bereits vorhandenen Communityportale durch das neue System zu ersetzen, welche gegenwärtig auf dem Framework \textbf{Drupal} 5 (PHP) basieren
 \item Eine moderne Suchfunktion durch einen Suchdaemon, z.B. Sphinx\footnote{\url{http://sphinxsearch.com/}} oder Lucene\footnote{\url{http://lucene.apache.org/java/docs/index.html}}
 \item Softwarestack: Ruby 1.9.2 mit Ruby on Rails 3.1, JavaScript\index{JavaScript} mit JQuery 1.6
 \item Die Software soll eine möglichst einfache und eingängige Bedienung haben
\end{itemize}


\section{Aufbau der Arbeit}


\paragraph{Hinweis zur Notation}


\subsection{Beispiel Syntax HL}


\begin{listing}
\begin{rubycode}
class Array
  # method documentation
  def mean
    self.inject(:+) / self.count.to_f
  end
end

"strings"
\end{rubycode}
\label{lst:example}
\caption{Ruby method definition with beautiful syntax hl}
\end{listing}


\newpage\chapter{Fazit}
\label{sec:fazit}
\epigraph{Write the tests you wish you had. If you don't, you will eventually break something while refactoring. Then you'll get bad feelings about refactoring and stop doing it so much. Then your design will deteriorate. You'll be fired. Your dog will leave you. You will stop paying attention to your nutrition. Your teeth will go bad. So, to keep your teeth healthy retroactively test before refactoring.}{Kent Beck}

\section{Zusammenfassung}


\newpage
 %\include{parts/01}
 %\section{Projektbeschreibung und Projektziele}
Für den praktischen Teil dieser Arbeit soll anhand der Entwicklung einer Ruby-on-Rails\hyp{}Anwendung die Testgetriebene\index{TDD} Entwicklung erprobt und angewendet werden.

Als Ergänzung zu den lokalspezifischen Communitys mit strenger Mitgliederauswahl soll nun ein neues, allgemeines IT-Jobportal entwickelt werden. Der voraussichtliche Name wird IT-Jobs\index{IT-Jobs-Projekt}-Und-Stellen.de\footnote{\url{http://www.it-jobs-und-stellen.de/}} sein.

Ziel soll es sein, den regionalen Organisationen auch eine Alternative zu den Community-Mitgliedschaften anzubieten, um kurzfristigen Personalbedarf zu decken, analog zu den bekannten Online-Stellenbörsen stepstone.de, monster.de und jobscout24\footnote{\url{http://www.stepstone.de} - \url{http://www.monster.de} - \url{http://www.jobscout24.de}}. Parallel zu dieser Arbeit wurden die Anforderungen analysiert und die in Abbildung \ref{fig:usecases} dargestellten Anwendungensfälle bestimmt.


\begin{figure}[htbp]
 \centering
 \includegraphics[width=1\textwidth]{./material/usecases.png}
 % usecases.png: 930x899 pixel, 51dpi, 46.50x44.95 cm, bb=
 \caption{Anwendungsfälle}
 \label{fig:usecases}
\end{figure}


\subsection{Anwendungsfälle}
Es gibt zwei verschiedene Nutzertypen. Zum einen der \textit{Bewerber}, der das Ziel hat, einen Job zu suchen, zu finden und sich darauf zu bewerben. Zum anderen der Mitarbeiter eines Kundenunternehmenes (\textit{Arbeitgeber}), der Stellenanzeigen auf der Webseite schalten möchte.
Verbal formuliert sind die Anwendungsfälle die Folgenden:
\begin{itemize}
 \item Ein Bewerber kann die Webseite nach sichtbaren Stellenangeboten durchsuchen.
 \item Ein Bewerber kann die Detailanzeige einer Stelle betrachten.
 \item Ein Bewerber kann sich auf eine Stelle über das System bewerben. Dabei kann er auch eine Verbindung zu seinem Facebook- oder LinkedIn-Account herstellen, falls er das möchte, um so automatisch einen Lebenslauf generieren zu lassen und Stammdaten eintragen zu lassen.
 \item Ein Kunde kann sich am Portal registrieren, dann anmelden und wird so zum \textit{Arbeitgeber}.
 \item Ein Arbeitgeber kann eine Stellenanzeige für ein Vielfaches von 30 Tagen schalten. Damit wird ein Bezahlvorgang über ein Drittanbieterportal (z.B. PayPal) ausgelöst und bei Bestätigung die Stelle sichtbar geschaltet. Er kann dazu zwischen 5 verschiedenen Designs wählen sowie Logo und Bild hinzufügen.
 \item Ein Arbeitgeber kann eine Stelle verlängern oder kopieren, um als sie als Vorlage für ein neues Stellenangebot zu nutzen.
 \item Ein Arbeitgeber kann eine laufende Stelle erneuern (refresh), um so seine Platzierung verbessern. Dies ist möglich, weil die Aktualität einer Stelle Auswirkungen auf die Platzierung in den Suchergebnissen hat. Ein Refresh kostet Geld und löst somit wieder einen Bezahlvorgang aus.
 \item Ein Arbeitgeber kann Statistiken über seine geschalteten Stellenanzeigen einsehen. Dies beinhaltet die Anzahl der Zugriffe pro Tag und Anzahl der Bewerbungen je Stelle.
 \item Ein Arbeitgeber kann seine Stellen mittels eines XML-Feed-Imports automatisiert einlesen. Dazu bezahlt er einen Pauschalbetrag je Kalenderjahr. Das System liest alle 6 Stunden automatisiert diesen Feed ein, der durch HTTP erreichbar sein muss und überträgt gültige Stellenanzeigen in die Datenbank\index{Datenbank}. Nicht mehr in dem Feed enthaltene Stellenanzeigen des Arbeitgebers werden in der Datenbank auf "`unsichtbar"' gesetzt. Das Format ist eine Modifikation von \glossar{RSS} 2.0.
\end{itemize}



\subsection{Nichtfunktionale Anforderungen}
Folgende Rahmenbedingungen und zusätzliche Ziele wurden vereinbart.

\begin{itemize}
 \item Eine hohe C0-\glossar{Testabdeckung} von mindestens 95\% als Grundlage für weitere Arbeiten und für den Prozess der Testgetriebenen Entwicklung\index{TDD}
 \item Eine hohe Erweiterbarkeit, um langfristig auch die bereits vorhandenen Communityportale durch das neue System zu ersetzen, welche gegenwärtig auf dem Framework \textbf{Drupal} 5 (PHP) basieren
 \item Eine moderne Suchfunktion durch einen Suchdaemon, z.B. Sphinx\footnote{\url{http://sphinxsearch.com/}} oder Lucene\footnote{\url{http://lucene.apache.org/java/docs/index.html}}
 \item Softwarestack: Ruby 1.9.2 mit Ruby on Rails 3.1, JavaScript\index{JavaScript} mit JQuery 1.6
 \item Die Software soll eine möglichst einfache und eingängige Bedienung haben
\end{itemize}


 %\include{parts/02}
 %\include{parts/03}

 %\input{parts/tools}
 %\input{parts/implementation}
 %\include{parts/04}

 \appendix

\chapter{Eigenschaften erfolgreicher Tests}

Das Vorhandensein von zahlreichen Tests reicht nicht aus, um die \glossar{quality} der Tests zu beuteilen. Stattdessen existieren einige Qualitätskriterien, um Tests auf ihre Sinnhaftigkeit und Wartbarkeit zu untersuchen:

\begin{description}
 \item[Unabhängigkeit (Independence) und Isolation] Ein Test \index{Test}ist unabhängig, falls er nicht durch andere Tests beeinflusst wird. Auch die Reihenfolge, in der die Tests ausgeführt werden, darf auf das Ergebnis keinen Einfluss üben. Siehe auch \citep{beck_test_2002}, \citep{palermo_guidelines_2006} sowie \citep[Karte 45]{langr_agile_2011}.
 \item[Wiederholbarkeit (Repeatability] Ein Test \index{Test}wird als wiederholbar bezeichnet, wenn er mehrmals hintereinander ausgeführt werden kann und dabei jedes Mal dasselbe Ergebnis liefert \citep[Karte 45]{langr_agile_2011} \citep{rappin_rails_2011}. Problematisch sind dabei z.B. Datum und Zeit sowie Zufallsfunktionen.
 \item[Klarheit (Clarity)] Ein Test \index{Test}ist klar, wenn sein Zweck sofort verständlich wird \citep{rappin_rails_2011}, \citep{palermo_guidelines_2006}. Damit wird einerseits die Lesbarkeit gemeint, anderseits schließt dies auch ein, ob der Test genau eine Eigenschaft testet und nicht redundant zu anderen Tests ist. Dies hat zur Folge, dass die Tests wartbarer werden und als Codedokumentation dienen können.
 \item[Präzise (Conciseness)] Ein Test \index{Test}ist präzise, wenn er so wenig Code und so wenige Objekte wie möglich benötigt, um sein Ziel zu erreichen \citep{palermo_guidelines_2006} \citep{rappin_rails_2011}. Eine Auswirkung ist, dass der Test schneller wird.
 \item[Robustheit (Robustness)] Ein Test \index{Test}ist robust, wenn es eine direkte Korrelation zum zu testenden Code gibt: Ist der Code korrekt, so ist der Test erfolgreich. Ist der Code falsch, so schlägt der Test fehl. Nicht-robuste Tests werden auch "`zerbrechlich"' (brittle) genannt. Dazu zählen auch sogenannte tautologische Tests, die immer erfolgreich verlaufen und keine Aussage über den zugrunde liegenden Programmcode geben.
 \item[Schnelligkeit] Gute Tests sollten schnell sein, da die Test-Suite\index{Test!Test-Suite} ansonsten deutlich seltener ausgeführt würd \citep{langr_agile_2011} \citep{palermo_guidelines_2006} \citep{nagappan_realizing_2008}. Siehe auch Präzision.
 \item[Zeitliche Nähe] Gute Tests entstehen in zeitlicher Nähe zu dem zugrunde liegenden Programmstück \citep{langr_agile_2011}. Innerhalb von TDD\index{TDD} ist dieses Kriterium immer erfüllt.
 \end{description}

% Microsoft Vorgehensweise:
%
% Start TDD from the beginning of projects. Do not stop in the middle and claim it
% doesn’t work. Do not start TDD late in the project cycle when the design has already
% been decided and majority of the code has been written. TDD is best done
% incrementally and continuously.
% For a team new to TDD, introduce automated build test integration towards the second
% third of the development phase—not too early but not too late. If this is a “Greenfield”
% project, adding the automated build test towards the second third of the development
% schedule allows the team to adjust to and become familiar with TDD. Prior to the
% automated build test integration, each developer should run all the test cases on their
% own machine.
% Convince the development team to add new tests every time a problem is found, no
% matter when the problem is found. By doing so, the unit test suites improve during the
% development and test phases.
% Get the test team involved and knowledgeable about the TDD approach. The test team
% should not accept new development release if the unit tests are failing.
% Hold a thorough review of an initial unit test plan, setting an ambitious goal of having
% the highest possible (agreed upon) code coverage targets.
% Constantly running the unit tests cases in a daily automatic build (or continuous
% integration); tests run should become the heartbeat of the system as well as a means to
% track progress of the development. This also gives a level of confidence to the team
% when new features are added.

% Encourage fast unit test execution and efficient unit test design. Test \index{Test}execution speed is
% very important since when all the tests are integrated, the complete execution can
% become quite long for a reasonably-sized project and when using constant test
% executions. Tests results are important early and often; they provide feedback on the
% current state of the system. Further, the faster the execution of the tests the more likely
% developers themselves will run the tests without waiting for the automated build tests
% results. Such constant execution of tests by developers may also result in faster unit
% tests additions and fixes.

% Share unit tests. Developers’ sharing their unit tests, as an essential practice of TDD,
% helps identify integration issues early on.

% Track the project using measurements. Count the number of test cases, code coverage,
% bugs found and fixed, source code count, test code count, and trend across time, to
% identify problems and to determine if TDD is working for you.

% Check morale of the team at the beginning and end of the project. Conduct periodical
% and informal surveys to gauge developers’ opinions on the TDD process and on their
% willingness to apply it in the future.

%
% % Einige dieser Punkte können durch Metriken überprüft werden. Dazu mehr im Abschnitt \ref{sec:metrics}.
% %  Code Craft S 144

\chapter{Nutzung von Cucumber in Verbindung mit Selenium für Firefox und Guard ohne X-Server}
Für ein effektives Test-Setup wurden folgende Testtools benutzt:
\begin{description}
 \item[Guard] automatische Testausführung bei Änderungen der Dateien.
 \item[Xvfb] Ein X-Server, der ohne Grafikausgabe arbeitet. Er eignet sich also für die Ausführung von GUI-Programmen wie Mozilla Firefox auf Remote-Servern.
 \item[Selenium\index{Selenium}] Interface zur Fernsteuerung von Browsern.
 \item[Spork] Erhöht die Testausführungsgeschwindigkeit, da Teile von Rails\index{Ruby-on-Rails} im Hintergrund gehalten werden und zwischen den Tests nicht neu geladen werden müssen.
\end{description}

\subsection*{Installation}
Es muss bereits installiert sein: firefox, xvfb, Rails > 3.0
\begin{enumerate}
 \item In das Gemfile folgendes eintragen:
 \begin{lstlisting}[caption=RAILS\_ROOT/Gemfile]
group :test, :cucumber do
  gem "capybara"
  gem 'cucumber'
  gem "cucumber-rails"
  gem 'database_cleaner'
  gem "launchy"
  gem "guard"
  gem 'guard-cucumber'
  gem 'guard-spork'
  gem "rb-inotify"  # Für Linux
  gem "spork", :git => "git://github.com/timcharper/spork.git"
end

 \end{lstlisting}
 Danach folgendes ausführen:

 \begin{verbatim}
bundle install
rails g cucumber:install
 \end{verbatim}


 \item Eine Datei Guardfile anlegen. Hier kommen alle Anweisungen zur Steuerung von guard hinein.
 \begin{lstlisting}[caption=RAILS\_ROOT/Guardfile]
group "cucumber" do
  guard 'spork' do
    watch('config/application.rb')
    watch('config/environment.rb')
    watch('features/support/env.rb')
    watch(%r{^config/environments/.+\.rb$})
    watch(%r{^config/initializers/.+\.rb$})
    watch('spec/spec_helper.rb')
  end
  guard 'cucumber', :cli => '--no-profile --color --format pretty --strict RAILS_ENV=test' do
    watch(%r{^features/.+\.feature$})
    watch(%r{^features/support/.+$})                      { 'features' }
    watch(%r{^features/step_definitions/(.+)_steps\.rb$}) { |m| Dir[File.join("**/#{m[1]}.feature")][0] || 'features' }
  end
end
\end{lstlisting}
\item In der Datei \texttt{config/cucumber.yml} die Option \texttt{--drb} für "`default"' und "`wip"' setzen.


\item Die Datei \texttt{features/support/env.rb} anpassen, um sie mit Spork kompatibel zu machen:
\begin{lstlisting}[caption=features/support/env.rb]
require 'cucumber/rails'
require "spork"
require 'test/unit/testresult'
Test::Unit.run = true

Spork.prefork do
  ENV["RAILS_ENV"] ||= "test"
  require File.expand_path(File.dirname(__FILE__) + '/../../config/environment')
  require 'cucumber/formatter/unicode'
  require 'cucumber/rails'
  require 'test/unit/testresult'
  ActionController::Base.allow_rescue = false
  begin
    DatabaseCleaner.strategy = :truncation
  rescue NameError
    raise "You need to add database_cleaner to your Gemfile (in the :test group) if you wish to use it."
  end
end

Spork.each_run do
  require 'cucumber/rails/world'
end

\end{lstlisting}

\end{enumerate}


\subsection*{Nutzung}

\begin{enumerate}
 \item Aktivierung von Xvfb auf Displayport 99
 \begin{verbatim}
  Xvfb :99 -ac
 \end{verbatim}

 \item Aktuelle Shell auf diesen Displayport einstellen
 \begin{verbatim}
  export DISPLAY=:99
 \end{verbatim}

 \item Guard starten
 \begin{verbatim}
  guard -g cucumber
 \end{verbatim}
\end{enumerate}

Nun werden automatisch alle Cucumber\index{Cucumber}-Features getestet. Falls Selenium\index{Selenium} verwendet wird, dann wird Firefox im Hintergrund gestartet, ohne dass dies sichtbar wird.






 \newpage
 \printglossary[toctitle=Glossar]
 %\setstretch{1,0} % Mein TOC ist zu lang!!
\listoffigures
%\setstretch{1,25} % Mein TOC ist zu lang!!

\addcontentsline{toc}{chapter}{Quellcode-Listings}
\lstlistoflistings

%  \cite{*} % Alles Zitieren DEBUG
\bibliographystyle{alphadin}     % Kurz
%\bibliographystyle{natbib}     % Verbos

\renewcommand{\indexname}{Stichwortverzeichnis}
%\addcontentsline{toc}{chapter}{Stichwortverzeichnis}
\printindex

\begin{raggedright}
  \bibliography{bibo}
\end{raggedright}

 \rowcolors{1}{white}{white}


\chapter*{Eidesstattliche Erklärung}
	Hiermit versichere ich, die vorliegende Arbeit selbstständig und unter ausschließlicher Verwendung
	der angegebenen Literatur und Hilfsmittel erstellt zu haben.

	Die Arbeit wurde bisher in gleicher oder ähnlicher Form keiner anderen Prüfungsbehörde vorgelegt
	und auch nicht veröffentlicht.\bigskip

Ich erkläre mich einverstanden, dass meine Diplomarbeit an Personen, die nicht mittelbar oder unmittelbar an meiner Prüfung beteiligt sind, ausgeliehen wird.\bigskip

\begin{tabular}{ll}

	Dresden, \today & \underline{\qquad \qquad \qquad \qquad \qquad \qquad}\\
	& \small{\ Unterschrift (Stefan Wienert)}
\end{tabular}

\rowcolors{1}{lightgray!50!white}{white}



\end{document}
